\subsection{Auxiliary topology and common localization}
\begin{lemma}[Pathwise finiteness from capped series]
For concrete zero-start decompositions $X_k=N_k+A_k$, set

\[
 B_k(t)=\sqrt{[N_k]_t}+\Var_{[0,t]}A_k.
\]
If $\sum_k E[1\wedge B_k(n)]<\infty$ for every integer $n\geq0$, then
$\sum_kB_k(t)<\infty$ for all times on one full-measure set.
\end{lemma}
\begin{proof}
The clocks are finite, nonnegative, increasing and measurable at each time.
At integer $n$, Tonelli gives almost-sure finiteness of the capped sum.
Its terms tend to zero, hence eventually $B_k(n)<1$ and the cap disappears.
The finitely many initial terms are finite too. Intersect over integers;
for any $t$ use monotonicity and an integer above $t$.
\end{proof}
This does not imply integrability of the uncapped sum.
\begin{lemma}[Monotonicity and right continuity of the sum]
On one full-measure set, $U_t=\sum_kB_k(t)$ is finite, increasing and
right-continuous.
\end{lemma}
\begin{proof}
Cumulative variation of a right-continuous locally BV path is right-continuous:
its right jump is the absolute right jump of the path, hence zero.
Each clock is therefore right-continuous. On $t\leq s<t+1$ it is dominated
by the summable $B_k(t+1)$. Dominated convergence for series gives right
continuity of the sum; monotonicity is termwise.
\end{proof}
\begin{lemma}[A common passage for concrete decompositions]
Under the usual conditions choose a strongly adapted c\`adl\`ag version of
$U$. Its strict absolute passage above $r+1$, capped at horizon $r+1$, defines
one localizing sequence $\tau_r\leq r+1$. On one full-measure set,

\[
 \sum_kB_k(t)\leq r+1
\]
 for every $r$ and $t<\tau_r$.
\end{lemma}
\begin{proof}
Cumulative variation at $t$ is recovered from a countable supremum of finite
grid variations before $t$, so is $\F_t$-measurable. The quadratic clocks
and their sum are adapted. Finite increasing paths have left limits;
right continuity gives c\`adl\`ag regularity. Correct the exceptional null
set at time zero. The c\`adl\`ag absolute-passage construction is increasing
and exhaustive. Before passage its magnitude is at most $r+1$;
indistinguishability returns the bound to the sum.
\end{proof}
The capped estimates are inputs here. The closed-graph comparison below will
supply them from Émery convergence.
\begin{lemma}[Boundary correction for a stopped quadratic root]
On one full-measure set,
\[
 \sqrt{[N]_t}\leq\sqrt{[N]_{t-}}+|\Delta N_t|
\]

for every $t$. For finite-valued stopping time $\tau$,
\[
 E\sup_{t\geq0}\sqrt{[N^\tau]_t}
 \leq E[\sqrt{[N]_{\tau-}}+|\Delta N_\tau|].
\]
A deterministic bound on $\tau$ and finite expectations are unnecessary.
\end{lemma}
\begin{proof}
Use $[N]_t=[N]_{t-}+(\Delta N_t)^2$ with nonnegative left limit.
The stopped increasing process attains its all-time supremum at $\tau$.
Evaluate the pathwise inequality there and integrate. At zero use the
left-limit convention and zero initial variation.
\end{proof}
\begin{lemma}[Prelocal cost from prefix components and jumps]
Under the usual conditions, for a zero-start decomposition $X=N+A$ and
finite stopping time $\tau\leq T$,
\[
 \mathcal J_{\tau-}(X)\leq
 E[\sqrt{[N]_{\tau-}}+|\Delta N_\tau|]+E[\Var(A)_{\tau-}].
\]
Local finite variation of $A$ suffices; expectations may be infinite.
\end{lemma}
\begin{proof}
First stop $A$ at the deterministic horizon $T$. This gives global bounded
variation without changing $A$ on $[0,\tau)$ or its strict prefix.
Monotonicity of cumulative variation gives
\[
 \Var(A)_t\leq\Var(A)_{\tau-}\quad(t<\tau).
\]
The strict-prefix variation estimate therefore yields
\[
 \Var(A^{\tau-})_\infty\leq\Var(A)_{\tau-}.
\]
The process $N^\tau+A^{\tau-}$ extends $X$ before $\tau$, with these two
components as an admissible decomposition. Bound the martingale cost by
the preceding quadratic-root boundary estimate, and integrate the global
variation bound for the other component. Finally take the defining infima
over extensions and decompositions.
\end{proof}
\begin{lemma}[Left-limit sums at the common passage]
For the same common passages constructed from the capped-series assumption,
on one full-measure set, for every $r$,
\[
 \sum_k\left(\sqrt{[N_k]_{\tau_r-}}+\Var(A_k)_{\tau_r-}\right)\leq r+1.
\]
\end{lemma}
\begin{proof}
The left limit of each clock is the sum of the square root of the left
limit of quadratic variation and the left limit of cumulative variation.
Fix a finite partial sum. For $t<\tau_r$ it is at most the full clock sum,
hence at most $r+1$. If $\tau_r>0$, take the left limit of this finite sum.
Then take the supremum over finite subsets to obtain the nonnegative-series
bound. If $\tau_r=0$, every left limit equals the zero initial value,
so the same conclusion holds.
\end{proof}
\begin{lemma}[Summable costs and witnesses at common stops]
The same $\tau_r$ satisfies
\[
 \sum_k\mathcal J_{\tau_r-}(X_k)\leq r+1+\sum_kE|\Delta N_{k,\tau_r}|.
\]
Define
\[
 s(N)=\sup_{\tau\text{ bounded stopping time}}E|\Delta N_\tau|.
\]
If $\sum_ks(N_k)<\infty$, exact prefix-representation witnesses can be
chosen at each $r$ with total cost at most

\[
 6(r+1+\sum_ks(N_k)+1).
\]
\end{lemma}
\begin{proof}
Stopped jumps are measurable. Sum the preceding single-decomposition cost
bounds over finite subsets; integrate the left-limit clock bound to get
$r+1$, then take the supremum over subsets. Each bounded $\tau_r$ makes the
jump term at most $s(N_k)$. Approximate the summable prelocal infima with
positive errors totaling one and apply the component selection bound with
Davis constant $6$.
\end{proof}
These are witnesses for arbitrary decompositions, not yet original-price
integral certificates.
\begin{definition}[The zero-start auxiliary cost]
For one decomposition $D=(N,A)$ and its quadratic variation, put
\[
 c_{r^1}(D)=s(N)+\sum_{n=0}^\infty2^{-n-1}
 E[1\wedge(\sqrt{[N]_n}+\Var(A)_n)],\qquad
 R(X)=\inf_{X\sim N+A}c_{r^1}(N,A).
\]
The two terms use the same decomposition before taking their joint infimum.
Infinite values are permitted.
\end{definition}
\begin{lemma}[Finiteness on all decomposable zero-start processes]
If $X$ has an admissible zero-start decomposition, $R(X)\leq3$.
Conversely $R(X)<\infty$ implies such a decomposition exists.
Finite global cost $\mathcal J(X)$ is not required.
\end{lemma}
\begin{proof}
Apply all-time DDY with threshold $1$ to the martingale component:
$N=L+Q$, with $|\Delta L|\leq2$. Then $X\sim L+(Q+A)$ is admissible.
Its bounded-stop jump cost is at most $2$; each capped expectation is at most
one and the weights total one. The joint cost is at most $3$.
Conversely the infimum over no decompositions is infinity.
\end{proof}
Regular zero-start good integrators enter this space by the decomposition
construction and elementary completeness proved later.
\begin{lemma}[Comparison with global cost]
$R(X)\leq2\mathcal J(X)$ as extended values. Thus vanishing global costs
imply vanishing auxiliary costs.
\end{lemma}
\begin{proof}
Fix a decomposition $D=(N,A)$ and denote its global cost by
\[
 j(D)=E\left[\sup_t\sqrt{[N]_t}+\Var(A)_\infty\right].
\]
The jump-square identity and monotonicity give, on one event for all times,
\[
 |\Delta N_\sigma|\leq\sqrt{[N]_\sigma}\leq\sup_t\sqrt{[N]_t}.
\]
Take expectations and the supremum over bounded stopping times to get
$s(N)\leq j(D)$. Every clock is bounded by the all-time root plus total
variation, so each capped expectation is also at most $j(D)$.
The geometric weights total one, hence $c_{r^1}(D)\leq2j(D)$.
Take the infimum over decompositions. All these inequalities hold for
extended values; no conversion of an infinite cost to a real number is needed.
Finally $0\leq R(X_k)\leq2\mathcal J(X_k)$ gives the convergence assertion.
\end{proof}
\begin{lemma}[Auxiliary triangle inequalities and differences]
Under the usual conditions,
\[
 R(X+Y)\leq R(X)+R(Y),\qquad R(-X)=R(X),\qquad
 R(X-Y)\leq R(X)+R(Y).
\]
In particular, $R(Z_k-X)\to0$ implies $R(Z_{k+1}-Z_k)\to0$.
\end{lemma}
\begin{proof}
Fix two admissible decompositions and add their components. Construct
quadratic variation of the martingale sum. The root triangle inequality
and subadditivity of variation bound its clock by the sum of the two clocks.
Integrate $1\wedge(a+b)\leq(1\wedge a)+(1\wedge b)$, multiply by the
positive weights, and sum to obtain subadditivity of the capped term.
For the jump term, integrate the jump triangle inequality at each bounded
stopping time and then take the supremum. Measurability of stopped jumps
permits separation of the nonnegative integrals.
Take infima over decompositions and quadratic variations, keeping both
cost coordinates attached to the same decomposition throughout.
Negation preserves quadratic variation, variation, and absolute jumps,
so it preserves $R$. The difference estimate follows by addition and
negation. Apply it to
$Z_{k+1}-Z_k=(Z_{k+1}-X)-(Z_k-X)$ for the convergence assertion.
\end{proof}
\begin{lemma}[Auxiliary cost under closed stopping]
For any stopping time $\tau$, $R(X^\tau)\leq R(X)$.
If $\sigma\leq\tau$, then $R(X^\sigma)\leq R(X^\tau)$.
Thus arbitrary termwise stopped versions of a vanishing-cost sequence also
have vanishing cost.
\end{lemma}
\begin{proof}
Stop both components. At any bounded stopping time the stopped jump equals
the original jump before or at $\tau$ and zero afterward, so its jump cost
cannot increase. Stopped quadratic variation and variation bound the clocks.
Take the joint infimum. Apply absorption for the second assertion and squeeze
for convergence.
\end{proof}
\begin{lemma}[Common localization from vanishing auxiliary costs]
If $R(X_k)\to0$, choose one strictly increasing subsequence $f$ and one
localizing sequence $\tau_r\leq r+1$. At every $r$, exact prefix witnesses
for $X_{f(k)}$ have total cost at most $6(r+4)$, and their costs tend to zero.
\end{lemma}
\begin{proof}
First suppose $\sum_kR(X_k)<\infty$. Approximate the infima with errors
totaling one, giving decompositions whose joint cost sum is at most
$\sum_kR(X_k)+1$. Their jump costs are summable. For each fixed integer
horizon the positive time weight implies summability of capped expectations.
Apply the common-localization theorem to these same decompositions.
The witness sum is at most $6(r+1+\sum_kR(X_k)+2)$.
For a general null sequence choose $R(X_{f(k)})<2^{-k-1}$, giving the bound
$6(r+4)$. Finite nonnegative sums imply termwise convergence to zero.
\end{proof}
The later topology comparison supplies this hypothesis from Émery convergence.
The global-cost comparison already supplies it when $\mathcal J(X_k)\to0$.
\begin{lemma}[Separation]
Indistinguishable processes have equal auxiliary costs. Under the usual conditions,
\[
 R(X)=0\ \Longleftrightarrow\ X\sim0,\qquad
 R(X-Y)=0\ \Longleftrightarrow\ X\sim Y.
\]
\end{lemma}
\begin{proof}
Transfer each decomposition and its quadratic variation to an
indistinguishable target without changing its cost; the infima agree.
If $R(X)=0$, the constant sequence $X_k=X$ has cost sum zero.
Choose near-minimizing decompositions and apply common localization.
For every $r$, $\sum_k\mathcal J_{\tau_r-}(X)<\infty$.
This is a constant nonnegative series, so $\mathcal J_{\tau_r-}(X)=0$.
For adapted right-continuous $X$ and fixed $t$, the random variable
$|X_t|1_{\{t<\tau_r\}}$ is measurable and has expectation zero by the
prefix maximal estimate. Intersect over $r$ and use exhaustion to obtain
$X_t=0$. Intersect next over a countable right-dense skeleton and extend
the equality to every time by right continuity.
For an arbitrary raw target, $R(X)<1$ supplies an admissible decomposition;
its component sum is a regular representative to which the argument applies.
The zero decomposition and zero quadratic variation prove the converse,
together with invariance under indistinguishability. Applying this equivalence
to differences proves the second assertion.
\end{proof}
\begin{lemma}[Summable costs imply pathwise uniform summability]
If $\sum_kR(X_k)<\infty$, then almost surely, for every finite $T$,

\[
 \sum_k\sup_{t\leq T}|X_k(t)|<\infty.
\]
A vanishing-cost sequence has one subsequence with this property.
\end{lemma}
\begin{proof}
For regular targets define $G_{k,r}$ as the supremum of $|X_k|$ over the
countable right-dense skeleton strictly before $\tau_r$.
It is measurable and $EG_{k,r}\leq6\mathcal J_{\tau_r-}(X_k)$.
Tonelli gives summability almost surely for all countably many $r$.
For any $T$ choose $r$ with $T<\tau_r$; right-continuous approximation by
skeleton points bounds the whole horizon supremum by $G_{k,r}$.
For raw targets use regular sums of near-minimizing decompositions and
countable indistinguishability. Extract a summable-cost subsequence for the
last assertion.
\end{proof}
\begin{lemma}[Auxiliary convergence implies ucp convergence]
If $R(X_k-X)\to0$, then for each finite $T$ and $\varepsilon>0$,

\[
 P\{\exists t\leq T:|X_k(t)-X(t)|>\varepsilon\}\to0.
\]
No extra measurability or path regularity of the raw targets is assumed.
\end{lemma}
\begin{proof}
For summable costs use regular decompositions and factorial envelopes to
obtain maxima that are measurable up to completion. The preceding lemma makes
them finite and tending to zero almost surely. Their real-valued versions
converge in probability; finiteness identifies their tails with the extended
maximal envelopes. From any subsequence of a vanishing-cost sequence extract
a summable one. The subsequence criterion gives convergence of the original
tail probabilities. Apply to differences; a strict exceedance at some time
is contained in the corresponding maximal tail event.
\end{proof}
Cost contraction also gives ucp convergence after arbitrarily selected
termwise stopping times.
\begin{proposition}[Metric quotient and Cauchy sequences]
Let $\mathcal D_0$ be the zero-start processes with admissible decompositions.
The finite pseudometric $d(X,Y)=R(X-Y)$ has separated quotient equal to the
indistinguishability quotient. Metric convergence is equivalent to vanishing
$R$-error and implies ucp convergence. Every Cauchy sequence has one strictly
increasing $f$ with

\[
 \sum_kR(X_{f(k+1)}-X_{f(k)})\leq1;
\] these same differences have common
localizers and exact witnesses with cost sum at most $6(r+4)$ at stage $r$.
\end{proposition}
\begin{proof}
Differences remain decomposable, so DDY gives finiteness. The auxiliary
algebra gives symmetry and triangle; separation identifies the quotient.
Its distance is precisely the representative cost.
For a Cauchy sequence choose tail bounds $b_n\to0$ and increasing $f$ with
$b_{f(k)}<2^{-k-1}$. Sum the adjacent-distance bounds, then apply the
summable-cost common-localization theorem to exactly these differences.
\end{proof}
\begin{lemma}[A locally uniform limit for the same subsequence]
If $\sum_kR(Y_{k+1}-Y_k)<\infty$, there is one process $Y$ such that,
on one full-measure set, $Y_k\to Y$ uniformly on every finite interval.
This applies to the same Cauchy subsequence carrying the common-localization
witnesses in the preceding proposition.
\end{lemma}
\begin{proof}
Summability of auxiliary costs gives, simultaneously for all finite $T$,
\[
 \sum_kq_k(T)<\infty,\qquad
 q_k(T)=\sup_{t\leq T}|Y_{k+1}(t)-Y_k(t)|.
\]
Each term is finite. At any fixed time, the series of distances between
successive values is bounded by this series. Completeness of the real
numbers gives a pointwise limit. Define $Y_t$ by that limit.
The summable uniform step bounds $q_k(T)$ upgrade pointwise convergence to
uniform convergence on $[0,T]$. The definition of $Y_t$ does not depend
on the horizon. The previously chosen subsequence has summable difference
costs, so this construction applies without a further extraction.
\end{proof}
This is initially a raw limit. We next retain global component series to
construct its decomposition and prove convergence of all cost coordinates.
\begin{lemma}[Martingale component sums at a fixed stop]
If true martingales $M_k$ satisfy $|M_k(t)|\leq g_k$ for all times with
measurable $g_k\geq0$ and $\sum_kEg_k<\infty$, then $M_t=\sum_kM_k(t)$
is a strongly adapted true martingale. This applies to the centered stopped
witness components
\[
 \sum_k(N_k-N_k(0))_{t\wedge\tau_r}
\]
 at each fixed stage.
\end{lemma}
\begin{proof}
Tonelli makes $G=\sum_kg_k$ integrable and finite almost surely.
The component series is absolutely convergent at all times on that event,
and its sum is measurable at every time. Finite sums are martingales dominated
by $G$, so the dominated martingale-limit lemma applies.
For witnesses the stopped centered components are true martingales, and the
finite-horizon envelope dominates them at all times because the stop is
bounded by that horizon. Centering increases cost by at most a factor of two,
preserving envelope summability.
\end{proof}
\begin{lemma}[Centered finite-variation series]
Zero-start paths $a_k$ with $\sum_k\Var(a_k)<\infty$ have a uniformly
convergent all-time sum $a$ of finite variation, and
\[
 \Var\left(a-\sum_{k<n}a_k\right)\longrightarrow0.
\]
At a fixed witness stage this applies almost surely to
$a_k=A_k^{\tau-}-A_k^{\tau-}(0)$.
\end{lemma}
\begin{proof}
Zero start bounds $|a_k(t)|$ by variation, giving uniform convergence of the
series. The finite-variation limit and variation-distance convergence follow
from summable variation of consecutive partial-sum differences.
For witnesses, centering does not change variation and the strict prefix
is constant after the stop, so its all-time variation is bounded by the
finite-horizon witness variation. Tonelli gives summability almost surely.
Centering removes initial constants that variation alone cannot control.
\end{proof}
We now keep one global decomposition per difference rather than choose
components independently at successive localization stages.
\begin{lemma}[Common estimates retaining global decompositions]
For a Cauchy sequence in the quotient, choose one strictly increasing $f$
and global decompositions of its actual neighboring differences,
$X_{f(k+1)}-X_{f(k)}\sim N_k+A_k$.
Their auxiliary costs sum to at most $2$, and one localizing sequence
$\tau_r\leq r+1$ satisfies
\[
 \sum_k\left(E\sqrt{[N_k]_{\tau_r}}+
 E\Var(A_k^{\tau_r-})\right)\leq r+3.
\]
The variation is over the entire time axis.
\end{lemma}
\begin{proof}
Choose the subsequence with adjacent auxiliary distances summing to at most
one. Approximating the infima by positive errors totaling one gives one
sequence of decompositions with total cost at most two.
Stop using the common clock of these same decompositions. The expected
sum of the pre-stop roots and cumulative variations is at most $r+1$.
For each term,
\[
 \sqrt{[N_k]_{\tau_r}}\leq
 \sqrt{[N_k]_{\tau_r-}}+|\Delta N_{k,\tau_r}|,\qquad
 \Var(A_k^{\tau_r-})\leq\Var_{[0,\tau_r)}(A_k).
\]
The boundary-jump expectations are bounded by the same decompositions'
costs of jumps at bounded stopping times, whose sum is at most two.
Pass from finite sums of nonnegative integrals to their supremum to obtain
the claim. Since stopped quadratic variation is increasing and constant
after the stop, its root at the stop is its all-time root supremum.
\end{proof}
\begin{lemma}[Sum of the same global martingale components]
Set $M_t=\sum_kN_k(t)$. Each $M^{\tau_r}$ is a true martingale, and $M$
is a zero-start local martingale.
\end{lemma}
\begin{proof}
For fixed $r$, Davis gives

\[
 \sum_kE\sup_t|N_k(t\wedge\tau_r)|\leq6\sum_kE\sqrt{[N_k]_{\tau_r}}<\infty.
\]
A measurable finite-horizon envelope dominates the stopped component for all
time. Its integrability upgrades the stopped component to a true martingale;
the dominated series lemma then applies. Evaluation at a stopped time
commutes with the pointwise sum, giving exactly $M^{\tau_r}$.
Zero initial values remove the localizing indicator, proving the local property.
\end{proof}
\begin{lemma}[A regular representative of the martingale sum]
The same $M$ is indistinguishable from an everywhere c\`adl\`ag, strongly
adapted, zero-start local martingale.
\end{lemma}
\begin{proof}
At every common stop the envelope expectations are summable; Tonelli gives
pathwise uniform convergence of the stopped finite sums. Uniform c\`adl\`ag
limits are c\`adl\`ag. Intersect over the stops. For any time $t$, exhaustion
places $[0,t+1)$ before one stop, transferring that regularity to the global
sum. Each time value is measurable. Replace nonregular paths by zero on the
time-zero measurable exceptional set and transfer the local-martingale property
by indistinguishability.
\end{proof}
\begin{lemma}[Sum of the same global finite-variation components]
For the same decomposition sequence set $A_t=\sum_kA_k(t)$.
On one full-measure set $A$ is locally of finite variation and, for every finite $T$,
\[
 \sum_{k<n}A_k\longrightarrow A\quad\text{uniformly on }[0,T],\qquad
 \Var_{[0,T]}\left(A-\sum_{k<n}A_k\right)\longrightarrow0.
\]
\end{lemma}
\begin{proof}
The capped-clock estimates give
$\sum_k\Var_{[0,T]}(A_k)<\infty$ for every $T$ on one full-measure set.
Fix $T$ and put $a_k(t)=A_k(t\wedge T)$. Its all-time variation is at most
$\Var_{[0,T]}(A_k)$ and $a_k(0)=0$.
The zero-start variation-series theorem shows that $\sum_ka_k$ has global
finite variation and that its partial sums converge both uniformly and
in all-time variation distance. On $[0,T]$ we have $a_k=A_k$, so both
convergences transfer back to the original components. Every compact
interval is contained in some $[0,T]$, proving local finite variation of $A$.
\end{proof}
\begin{lemma}[A regular representative of the variation sum]
The preceding sum has an indistinguishable zero-start strongly adapted
representative with everywhere c\`adl\`ag locally BV paths.
\end{lemma}
\begin{proof}
Measurability of each series value gives adaptation. Local uniform convergence
of right-continuous finite sums gives right continuity of the limit.
Correct the common exceptional set for right continuity and local BV using
the usual conditions at time zero. Zero start is preserved, and local BV
gives left limits.
\end{proof}
The global components were fixed once: equality of sums would not by itself
identify separate components in decompositions with merely adapted FV parts.
\begin{lemma}[A limit inside the decomposable-process space]
Every auxiliary-Cauchy sequence has a subsequence $X_{f(n)}$ converging almost
surely locally uniformly to an admissibly decomposable process $Y$.
Its pointwise limit is indistinguishable from this $Y$ and also decomposable.
\end{lemma}
\begin{proof}
Use the same regular sums $M,A$ and set $Y=X_{f(0)}+M+A$.
Add the base term's decomposition to those of the sums.
On the common event of all component identities and convergence, choose
$\tau_r>T$ to transfer stopped martingale uniform convergence to $[0,T]$.
The FV sums also converge uniformly there. Telescoping gives

\[
 X_{f(0)}+\sum_{k<n}(N_k+A_k)=X_{f(n)},
\] hence convergence to $Y$.
Uniqueness of each time limit on that event gives indistinguishability and
transfers the decomposition.
\end{proof}
\begin{lemma}[Stopped-jump lower semicontinuity]
If martingale components $U_n$ tend almost surely locally uniformly to a
process $U$ with left limits, then $s(U)\leq\liminf_ns(U_n)$.
In particular the same regular sum satisfies
\[
 s(M)\leq\sum_ks(N_k).
\]
\end{lemma}
\begin{proof}
Fix bounded $\sigma\leq T$. Uniform convergence after deterministic stopping
at $T$ gives convergence of left limits and hence of jumps at $\sigma$.
Fatou yields $E|\Delta U_\sigma|\leq\liminf_ns(U_n)$, with bound independent
of $\sigma$; take the supremum. Apply to finite sums, whose jump cost is at
most the sum of their costs, using the already proved local uniform convergence.
\end{proof}
\begin{lemma}[Component-series tails]
For the same regular sums $M,A$ and original decompositions, for every $n$,
\[
 s\left(M-\sum_{k<n}N_k\right)\leq\sum_{k=0}^{\infty}s(N_{k+n}).
\]
The left side tends to zero. At every finite horizon $T$ we also have
\[
 E\left[\min\left\{\Var_{[0,T]}
       \left(A-\sum_{k<n}A_k\right),1\right\}\right]\longrightarrow0.
\]
\end{lemma}
\begin{proof}
Fix $n$ and shift the original decomposition sequence by $n$.
Its martingale partial sums converge locally uniformly to
$M-\sum_{k<n}N_k$ on the same full-measure set.
Apply stopped-jump lower semicontinuity to this sequence for the first
inequality. Its right side tends to zero since $\sum_ks(N_k)<\infty$.
For finite variation, use the pathwise variation convergence already proved.
The regular sum and every partial sum are strongly adapted and
right-continuous, so the variation of their difference on a finite interval
is measurable. Transfer to the same representatives by indistinguishability
and apply dominated convergence with the constant bound one.
\end{proof}
\begin{lemma}[Quadratic-root tails of the same martingale sum]
Let $U_n=M-\sum_{k<n}N_k$. At a common bounded stop $\tau$ with summable
original root expectations,
\[
 E\sqrt{[U_n]_\tau}\leq42\sum_{k\geq n}E\sqrt{[N_k]_\tau}\longrightarrow0.
\]
The intrinsic quadratic variation of each tail is independent of this stop.
\end{lemma}
\begin{proof}
If partial sums converge uniformly on $[0,T]$ to local martingale $U$, then

\[
 U_T^*\leq\sum_k(N_k)_T^*,
\] allowing an infinite envelope sum.
Reverse Davis, Tonelli, and forward Davis give

\[
 E\sqrt{[U]_T}\leq7EU_T^*\leq7\sum_kE(N_k)_T^*
\leq42\sum_kE\sqrt{[N_k]_T}.
\]
Apply this to the stopped components and their tail series. Their summable
envelope expectations give all-time uniform convergence almost surely to
$U_n^\tau$. Choose deterministic $T\geq\tau$ and use stopping consistency.
The right side is the tail of a summable nonnegative series.
\end{proof}
\begin{lemma}[Removing common stops from capped roots]
Suppose $\tau_r\to\infty$ almost surely and
$E\sqrt{[U_n]_{\tau_r}}\to0$ for every $r$. Then for each finite $T$,
\[
 E[\min\{\sqrt{[U_n]_T},1\}]\longrightarrow0.
\]
\end{lemma}
\begin{proof}
Monotonicity of quadratic variation gives
\[
 \min\{\sqrt{[U_n]_T},1\}
 \leq\sqrt{[U_n]_{\tau_r}}+1_{\{\tau_r\leq T\}}.
\]
The second term's expectation tends to zero as $r\to\infty$ by exhaustion
and dominated convergence. First fix $r$ making this error small, then
send $n\to\infty$ in the first term.
\end{proof}
\begin{theorem}[Completeness of the auxiliary metric]
Under the usual conditions, the indistinguishability quotient of all
zero-start processes admitting a decomposition is complete for
$d(X,Y)=R(X-Y)$.
\end{theorem}
\begin{proof}
Choose a fast subsequence $X_{f(k)}$ of a Cauchy sequence in the quotient,
retaining decompositions $(N_k,A_k)$ of its actual neighboring differences.
Let $M,A$ be the regular component sums constructed above and put
$X=X_{f(0)}+M+A$. This process admits a decomposition.
Subtracting the first $n$ terms from these same sums decomposes
$X-X_{f(n)}$ into
\[
 U_n=M-\sum_{k<n}N_k,\qquad B_n=A-\sum_{k<n}A_k.
\]
The identification of the target is the telescoping identity for the
original adjacent differences. The jump-tail estimate gives $s(U_n)\to0$,
and at each finite $T$,
\[
 E\min\{\sqrt{[U_n]_T},1\}\to0,\qquad
 E\min\{\Var_{[0,T]}(B_n),1\}\to0.
\]
The inequality $\min\{a+b,1\}\leq\min\{a,1\}+\min\{b,1\}$ therefore
makes the capped clock expectation of this same decomposition tend to
zero at every integer horizon. These expectations are at most one, so
dominated convergence also applies to their geometric weighted sum.
The decomposition's $r^1$ cost tends to zero, and hence so does its
infimum $R(X-X_{f(n)})$. Negation invariance gives convergence of the
subsequence to $X$ in the quotient. A Cauchy sequence with a convergent
subsequence converges entirely to the same limit. Any quotient sequence
can be lifted to representatives for this argument, proving completeness.
\end{proof}
\begin{lemma}[Continuous scalar multiplication]
For every decomposable zero-start process $X$,
\[R(cX)\leq\max(1,|c|)R(X).\]
If $a_k\to a$ and $R(X_k-X)\to0$, then $R(a_kX_k-aX)\to0$.
\end{lemma}
\begin{proof}
Scale one decomposition. The jump and square-residual characterizations give
$[cN]=c^2[N]$, so clocks and jump costs scale by $|c|$.
Use $1\wedge(|c|b)\leq\max(1,|c|)(1\wedge b)$ and take joint infima.
For a fixed $X$ and $a_k\to0$, that inequality alone is insufficient.
Choose a DDY decomposition with finite jump cost. Its scaled jump cost
vanishes; each finite clock is finite pathwise, so its scaled capped
expectation vanishes by bounded convergence. Apply dominated convergence
also to the geometric series. Hence $R(a_kX)\to0$.
Now split
\[
 a_kX_k-aX=a_k(X_k-X)+(a_k-a)X
\]
 and use boundedness of $(a_k)$,
the triangle inequality, and the fixed-process conclusion. Sequential
continuity in these metric spaces gives continuity, descending to the quotient.
\end{proof}
